\documentclass[english]{article}
\usepackage[utf8x]{inputenc}
\usepackage[T1]{fontenc}
\usepackage{babel}
\usepackage{amsmath}
\usepackage{graphicx}
\usepackage{hyperref}
\usepackage{fancyhdr}
\usepackage{booktabs}
\usepackage{lipsum}
\usepackage{cite}
\usepackage{setspace}
\usepackage{float}
\usepackage{verbatim}
\usepackage{caption}

\graphicspath{{../images/}}
\pagestyle{fancy}
\fancyhf{}
\renewcommand{\headrulewidth}{0pt}
\setlength{\headheight}{40pt} 
\usepackage[a4paper,top=1.75cm,bottom=2cm,left=2cm,right=2cm,marginparwidth=1.75cm]{geometry}

\begin{document}

%-------------------------------------------------------%

\title{\bf Designing Sustainable Agriculture Databases: Relational and Document-Based Approaches}
\author{Student Number - 720017170}
\date{\today}
\maketitle
\thispagestyle{fancy}

\vspace{-3em}
%-------------------------------------------------------%
\section{Introduction}

\vspace{-0.5em}

This report describes the design, implementation, and analysis of a data system which
stores UK agricultural sustainability initiatives. 
An efficient and reliable database is developed to store and analyse 
the performance of sustainability initiatives across UK farms.

\vspace{0.25em}

\textbf{The project discusses four key sections:}

\vspace{-0.5em}

\begin{itemize}
\item \textbf{Relational Database Design:} A relational database design adhering to
the third normal form ($\text{3NF}$), as well as an entity relationship diagram ($\text{ERD}$),
supported by detailed design rationale.

\vspace{-0.4em}

\item \textbf{Database Implementation:} An outline and \texttt{.sql} script, implemented in MySQL.
The script is used to create and populate the database tables from raw, unnormalized data.

\vspace{-0.4em}

\item \textbf{RESTful API Design \& Implementation:} The design and implementation
of a RESTful API that allows for data access and modification within the database,
including all endpoints and their functions.

\vspace{-0.4em}

\item \textbf{Alternative Design:} An alternative design using a document-based NoSQL
database model and a comparison with a relational approach, considering the trade-offs in scalability, 
data integrity, flexibility, and performance.

\end{itemize}

\vspace{-2em}

%-------------------------------------------------------%
\section{Database Design}

\vspace{-0.75em}

The database was designed to 3NF, where the key considerations were data integrity, reducing redundancy, 
and ensuring efficient data retrieval; all of which are important for designing
production-level databases\cite{Connolly2014-tc}.

\vspace{-1em}

\subsection{Design Rationale}

\vspace{-0.5em}

The aim of the design was to remove all data redundancy and inconsistencies. 
This was achieved by ensuring all non-key attributes are fully dependent only on the
primary key of that table with no transitive dependencies\cite{Connolly2014-tc}. By adhering to 3NF, the design
will eliminate update, insertion, and deletion errors, which in turn increasing consistency and 
accuracy\cite{Connolly2014-tc}.

\vspace{-1em}

\subsection{Normalisation and Referential Integrity}

\vspace{-0.5em}

The design process began by considering first and second normal forms (1NF, 2NF). 
The raw data was processed to ensure 1NF by guaranteeing that each row and column contains 
only a single value, and all duplicate groups were removed\cite{Connolly2014-tc}. To achieve 2NF, 
each non-primary key attribute was made fully dependent on only the primary key of its table\cite{Connolly2014-tc}.
For example, attributes such as crop\_name and resource\_type were dependent only on 
their tables primary keys, not on any other records. Consequently, descriptive tables (Crop, Resource), 
which stored each value only once, were created to eliminate repeated data storage\cite{Connolly2014-tc}.

Several junction tables were created to implement 3NF, as all transitive
dependencies must be removed\cite{Connolly2014-tc}. For instance, the Farm\_Crop table resolves the many-to-many relationship 
between farms and crops, ensuring temporal attributes (planting\_date, harvest\_date) that are 
specific to each farm's growing of a crop, are accurately recorded. The Crop\_Resource table captures
the relationship between crops and resources, storing when each resource was applied to a crop and in what quantity. 

To enable the analysis of initiative impact on crop performance,
Crop\_Initiative was designed to link crops to their respective sustainability initiative. 
Also, the Farm\_Initiative table consolidates the relationship between farms and 
initiatives, storing initiative-specific metrics for each farm (date\_initiated, 
expected\_impact, ev\_score, water\_source). This allows the same initiative to be tracked 
separately across different farms, with potentially different values per farm\cite{Connolly2014-tc}.

Labour hours are captured in the Labour table using the unique combination of farm, crop, resource,
and initiative. The combination is necessary because labour requirements are dependent on which 
resources are applied to a specific crop on a farm that is using a particular initiative.

Finally, all junction tables use surrogate primary keys for efficient 
indexing, while foreign key constraints enforce referential integrity and prevent data 
inconsistency across the database.

\vspace{-1em}

\subsection{Entity Relationship Diagram (ERD)}

\vspace{-1em}

\begin{figure}[H]
    \includegraphics[width=1\textwidth]{../images/E-R.pdf}
    \vspace{-1.75em}
    \caption{Normalised (3NF) Entity Relationship Diagram}
    \label{fig:erd}
\end{figure}

\vspace{-1.5em}
%-------------------------------------------------------%
\section{SQL Implementation}

\vspace{-0.5em}

Section 3 discusses the processes taken to develop the database within MySQL from the E-R diagram
in Figure~\ref{fig:erd}, with more detailed implementation and SQL provided in the accompanying 
\texttt{.sql} script.

\vspace{-1em}

\subsection{Database Configuration}

\vspace{-0.5em}

First, a new database environment was created within MySQL Workbench, making sure the database implementation is clean 
and reproducible. The process removes any previous version of the database in that environment. Following this, 
the database can be created using the \texttt{CREATE DATABASE} command. To allow file population from the 
CSV provided, local file loading is enabled. The specific set up guarantees that the database is reproducible from a 
blank baseline so that all future commands work within the environment.

\vspace{-1em}

\subsection{Table Creation}

\vspace{-0.5em}

After the database was created and configured, tables were created to store each record.
Initially, descriptive tables were created to contain data that only related to Farms, Soils, Crops, 
Initiatives and Resources. Junction tables also were defined in order to facilitate many-to-many 
relationships as well as store instance specific attributes. These included Farm\_Crop, Crop\_Resource,
Crop\_Initiative, Farm\_Initiative, and Labour. Each table was created with appropriate primary keys, foreign keys 
and constraints to ensure data integrity. 

\vspace{-1em}

\subsection{Data Loading and Cleaning}

\vspace{-0.5em}

First, the raw data provided was loaded into a staging table using the MySQL \texttt{LOAD DATA INFILE} command.
The staging table mirrors the structure of the raw CSV file, ensuring a clean import, as well as providing a 
location for data cleaning. Once within the staging table, certain columns required cleaning
to ensure compatibility with the final database schema. Cleaning included standardising date formats, due to 
formatting inconsistencies resulting from loading.

After cleaning, the data was then inserted into the schema using \texttt{INSERT INTO ... SELECT} statements.
These statements extract the relevant data from the staging table, and load it into
the final schema. The cleaning process meant that all data in the final tables adhered to the defined structure
and integrity constraints of the database schema. After population, the database was tested with various queries to 
validate the design and implementation, all shown at the end of the \texttt{.sql} script.

\vspace{-1em}

\subsection{Cloud Extension}

\vspace{-0.5em}

Beyond the core requirements, the database was set up in a cloud environment using Microsoft Azure's SQL Database 
service. This involved creating a new SQL schema within Azure using MSSQL. The existing \texttt{.sql} script was adapted 
to be compatible with Azure's SQL syntax and features. A Python script was developed to automate the data loading
process into the Azure database, ensuring that the data could be efficiently and accurately populated
in the cloud environment. The cloud implementation facilitates RESTful API access.

\vspace{-1em}

%-------------------------------------------------------%
\section{RESTful API Design}

\vspace{-0.5em}

\subsection{Design}

\vspace{-0.5em}

To allow the database to be accessible from other applications, a RESTful API was designed to facilitate 
data access and modification. The API uses standard HTTP methods (GET, POST, PUT, PATCH, DELETE) which perform
CRUD operations within the database\cite{SelvarajSivaraj2024ItA}.

The design utilises generic endpoints, shown in Table~\ref{tab:api-endpoints}, able to access and manipulate data across
all tables. For example, \texttt{GET /:table/:id} gives access to all data from
any table by specifying the table name and an optional record ID, implementing request flexibility along with code 
reusability and simplicity during implementation.
The generic approach does sacrifice some specificity and optimisation that more tailored endpoints may achieve. 
The \texttt{/query} endpoint permits custom SQL queries, providing flexibility to data analysts.
Flexibility reduces constraints caused by predefined endpoints. Letting users enter custom SQL could 
introduce security risks if not properly managed; however, security considerations can be mitigated by 
restricting access to read-only queries and implementing access control.

In a production environment, additional tailored endpoints could be implemented for frequently accessed 
resources, such as \texttt{GET /farms/{farmID}/environmental-score} to retrieve sustainability metrics 
for specific farms, or \texttt{GET /crops/{cropID}/yield} to access crop performance data. 
These would optimise performance for common use cases while the generic endpoints provide flexibility for 
ad-hoc analysis.

All data is exchanged in JSON format, which is widely used and supported across various programming 
languages and platforms\cite{SelvarajSivaraj2024ItA}. The design along with concise error codes and messages, 
makes the API user-friendly and accessible for developers interacting with the agricultural database.

\vspace{-1em}

\renewcommand{\arraystretch}{1.6} 
\begin{table}[H]
\centering
\caption{RESTful API Endpoint Definitions}
\vspace{-1em}
\label{tab:api-endpoints}
\scriptsize	
\begin{tabular}{@{}llp{5cm}p{3.5cm}p{4cm}@{}} 
\toprule
\textbf{Method} & \textbf{Endpoint} & \textbf{Parameters} & \textbf{Description} & \textbf{Response} \\
\midrule

GET & \texttt{/:table} & 
\textbf{Query:} \texttt{limit} (int), \texttt{offset} (int), \texttt{sort\_by} (string), \texttt{order} (asc|desc) & 
Retrieves all records from specified table with pagination and sorting. & 
200 OK: JSON array \\

GET & \texttt{/:table/:id} & 
\textbf{Path:} \texttt{:id} (int) — primary key & 
Retrieves single record by primary key. Returns empty if not found. & 
200 OK: JSON object | 404 Not Found \\

POST & \texttt{/:table} & 
\textbf{Body:} JSON object matching table schema. Example for Farm: \texttt{\{"farm\_location": "string"\}} & 
Creates new record with provided column-value pairs. & 
201 Created: JSON \texttt{\{ok, rowsAffected\}} | 400 Bad Request \\

PUT & \texttt{/:table/:id} & 
\textbf{Path:} \texttt{:id} (int) \newline
\textbf{Body:} Complete JSON object with all table fields & 
Full record replacement. All fields must be provided. & 
200 OK | 404 Not Found | 400 Bad Request \\

PATCH & \texttt{/:table/:id} & 
\textbf{Path:} \texttt{:id} (int) \newline
\textbf{Body:} Partial JSON object with only fields to update & 
Partial update. Only specified fields modified. & 
200 OK | 404 Not Found | 400 Bad Request \\

DELETE & \texttt{/:table/:id} & 
\textbf{Path:} \texttt{:id} (int) & 
Deletes record by primary key. & 
200 OK: \texttt{\{ok, message\}} | 404 Not Found \\

POST & \texttt{/query} & 
\textbf{Body:} \texttt{\{"sql": "SELECT ... FROM ..."\}} \newline
\textbf{Restriction:} SELECT-only queries permitted & 
Executes read-only SELECT queries, others blocked for security. & 
200 OK: JSON array | 400 Bad Request | 403 Forbidden \\

\bottomrule
\end{tabular}
\end{table}
\renewcommand{\arraystretch}{1}  % Reset to normal after table




\vspace{-2.5em}

\subsection{Implementation Extension}

\vspace{-0.5em}

To extend the coursework and validate my design, I developed the RESTful API using Node.js, Express and MSSQL to
access and query the cloud database. The API implements all the designed endpoints outlined in Section 4.1, and is
hosted on Azure App Service, providing a scalable and reliable platform for deployment.
The deployed API can be accessed at: 
\url{http://agriculture-api-jlg-g2gffahdgxfqc6b3.swedencentral-01.azurewebsites.net/} and the interface
is available at: \url{https://joshlg18.github.io/Agricultural-Database/}.

\vspace{-1em}

%-------------------------------------------------------%
\section{Document Based Design Implementation}

\vspace{-0.5em}

\subsection{Limitations of Relational Models}

\vspace{-0.5em}

The relational database model maintains data integrity and structured querying capabilities; however
it can introduce issues in scenarios requiring high flexibility and scalability\cite{Connolly2014-tc}.

Relational databases struggle with handling unstructured or semi-structured data\cite{alma991016914362807446}, which is 
becoming more common especially in the agricultural context where satellite images and sensor data may be involved.
The rigid nature of relational schemas makes it difficult to adapt to evolving data
requirements without significant schema restructuring\cite{Connolly2014-tc}.

The limitations described emphasise the need to explore alternative database designs such as NoSQL document-based models,
which are built for greater flexibility and scalability\cite{alma991016914362807446}.
  
\vspace{-0.5em}

\subsection{Alternative NoSQL Design}

\vspace{-0.5em}

To overcome the constraints of the relational model, a document-based NoSQL database design could be introduced.
NoSQL approaches focus on flexibility and scalability\cite{alma991016914362807446}, making it well suited for handling unstructured
and semi-structured data. Within this database, collections would contain key entities relating
to farms, crops, resources, and initiatives, stored in JSON or BSON format, supporting nested structures and arrays.

Instead of relying on joins between tables through relationships, documents would embed related data within them.
For example, a farm document could contain an array of crop documents, each with its own set of resources and 
initiatives.  The denormalised approach reduces the need for complex joins, improving read performance and simplifying 
data retrieval\cite{alma991016914362807446}, but can have some limitations presented in Section 5.3. 
Example collections are shown in Figure~\ref{fig:nosql-docs}.

\vspace{-0.5em}

\begin{figure}[H]
\centering
\footnotesize

\begin{minipage}{0.5\linewidth}
\begin{verbatim}
Farm data is stored within a farm collection as:
{
  "farmID": "1",
  "farm_location": "South Farm, Kent",
  "crops": [
    {
      "cropID": "101",
      "crop_name": "Wheat",
      "resources": [{ "resourceID": "1" }],
      "initiatives": [{ "initiativeID": "1" }]
    }
  ]
}
Initiative data is stored within an initiative collection as:
{
  "initiativeID": "1",
  "description": "Organic Farming"
}
\end{verbatim}
\end{minipage}

\vspace{-1em}

\caption{Example NoSQL Document Structure for Farm and Initiative Collections}
\label{fig:nosql-docs}
\end{figure}


\vspace{-2.5em}

\subsection{Comparative Analysis}

\vspace{-0.5em}

While the NoSQL document-based design provides improvements in flexibility and scalability, it can introduce challenges
compared to the relational approach. The relational model relies on vertical scaling, by enhancing the capacity of a 
single server, which could be more straightforward for smaller datasets and simpler applications. In contrast, 
NoSQL databases are horizontally scalable, with the ability to distribute data across multiple 
servers\cite{alma991016914362807446}. The NoSQL implementation enables nearly
constant growth, making it suitable for deployment in the event the scale of the agricultural system expands. 

However, the denormalised structure of NoSQL databases can lead to data redundancy and potential inconsistencies, 
as related data can be duplicated across multiple documents\cite{alma991016914362807446}. Duplication can complicate 
data management, especially when updates are required. For example, if a farm changed name or an initiative altered 
its description, that data will need to be altered across multiple documents. In contrast, within relational databases 
these changes can occur in a single location, the Farm or Initiative table respectively, improving 
consistency\cite{Connolly2014-tc}.

Schema flexibility presents another important consideration between the two approaches. Document-based databases 
can have dynamic schemas, meaning that different documents can have different data types and 
structures\cite{alma991016914362807446}. Conversely, relational databases require a fixed schema, which provides structure and consistency 
but can be less adaptable to alterations\cite{Connolly2014-tc}, which may be important if data requirements shift e.g.
a farm needed to store unstructured data or additional attributes.

Performance also needs to be considered when weighing up the two approaches. 
Relational databases can excel in complex queries and transactions due to their structured nature and support 
for SQL\cite{Connolly2014-tc}. NoSQL databases, on the other hand, provide superior performance for read-heavy 
workloads and large-scale data retrievals, as they are designed to handle high volumes of unstructured 
data efficiently\cite{alma991016914362807446}.

The choice between approaches can be dependent on the data and project requirements. For the current agricultural 
system, the relational model is more suitable because the data has stable, well-defined structures.
The system's primary use case involves complex analytical queries across multiple entities, which relational
databases handle efficiently. However, if the system were to migrate towards handling high-volume, potentially 
unstructured data e.g. sensors or images, transitioning to a NoSQL approach may become 
advantageous.

\vspace{-0.5em}

%-------------------------------------------------------%

% Include your references here.
\begin{spacing}{0.5}
    \bibliographystyle{IEEEtran}
    \bibliography{../references/references}
\end{spacing}

\vspace{-1em}

%-------------------------------------------------------%
\end{document}
